% Section 3.1, from lectures/L03/appendix/a_the_kernel_and_its_tree.md.

\section{The Kernel as a Program, and Its Source Tree}
\label{c3:app:a}

This section is the first half of the chapter: what kind of program the kernel is, what its one
public interface is and what follows from that, and how twenty-four million lines of C are arranged
so that you can find things in them. The configuration half is \secref{c3:app:b}.

\subsection{What kind of program the kernel is}
\label{c3:app:a1}

You have written programs. The kernel is a program, and almost every assumption you carry from the
first to the second is wrong. Five differences matter.

\textbf{It has no \code{main} and never returns.} It has an entry point that the bootloader jumps
to, it brings itself up, and then it stops being a thing that is ``running'' at all. After boot the
kernel is a body of code that gets entered: by a system call, by an interrupt, by a fault. Between
those it is not executing. Chapter~\ref{c4:ch}'s modules are the same shape, which is why
\code{module_init} returning does not mean the module is finished.

\textbf{It is one address space, shared by everything in it.} There is no isolation between
subsystems. A null pointer dereference in a sound driver corrupts the same address space the
filesystem code is using. This is what ``monolithic'' means, and it is worth being precise:
\textbf{monolithic is a claim about the address space, not about modularity.} The kernel is
intensely modular in its source organisation; it just does not enforce any of those boundaries at
run time.

\textbf{It cannot fault the way a process can.} A userspace segfault kills one process and the
system carries on, because the kernel is there to clean up. When the kernel dereferences a bad
pointer there is nobody above it. What happens instead is an \textbf{oops}: the kernel prints a
register dump and a backtrace and tries to kill the offending task. If it cannot, because the fault
was in something no task owns, it panics.

\textbf{It has no C library.} Compiled \code{-nostdinc}, \code{-ffreestanding}, against its own
headers. There is no \code{printf}, no \code{malloc}, and no floating point. Chapter~\ref{c4:ch}
covers the consequences; the point here is that ``the toolchain'' and ``the kernel'' are two of
Chapter~\ref{c1:ch}'s four pieces and not one, and this is why.

\textbf{Its stack is small and fixed.} 16 KB on arm64, for the entire call chain, and nothing grows
it. Run off the end and the best you can hope for is a guard page and a panic.

\subsection{The one public interface, and the one that is not}
\label{c3:app:a2}

This is the most consequential fact in the chapter, and it is two facts that people routinely merge
into one wrong one.

\textbf{The system call interface is stable forever.} A binary from 1998 still runs. Linus enforces this with unusual ferocity, and the rule is short: \emph{we do not break
userspace.} If a change makes a working program stop working, the change is reverted, however
correct it was.

\textbf{The internal API has no stability guarantee whatsoever.} Function signatures change,
structures gain and lose fields, whole subsystems are replaced. Between two point releases, code
inside the kernel may be rearranged freely.

Both are deliberate, and together they explain almost everything about how kernel development works:
\begin{itemize}
  \item \textbf{An in-tree driver is maintained for you.} When somebody changes an internal API they
    update every caller in the tree, including your driver, because the tree must always build.
  \item \textbf{An out-of-tree driver is maintained by you, forever.} Nobody changing an API knows
    your driver exists. This is the real cost of staying out of tree, and it is a much larger cost
    than the licensing argument in Chapter~\ref{c1:ch}. \code{vermagic} and
    \code{CONFIG_MODVERSIONS} in Chapter~\ref{c4:ch} are the mechanisms that turn the resulting
    mismatch into an error message rather than into corruption.
  \item \textbf{``Why is there no stable driver ABI'' is a settled question}, and the documented
    answer is in \file{Documentation/process/stable-api-nonsense.rst}. The short version is that a
    stable internal ABI would freeze design mistakes permanently, and the project has decided it
    would rather move the callers.
\end{itemize}

\subsection{The subsystem map}
\label{c3:app:a3}

Roughly, the kernel is:

\begin{booktable}{@{}>{\hsize=0.55\hsize}L>{\hsize=1.45\hsize}Ll@{}}
  \thead{Subsystem} & \thead{Responsible for} & \thead{Source} \\
  \midrule
  Process management & Creating, scheduling and destroying tasks & \code{kernel/} \\
  Memory management & Address spaces, page tables, allocation & \code{mm/} \\
  Virtual filesystem & The common file interface over many filesystems & \code{fs/} \\
  Networking & Sockets, protocols, packet handling & \code{net/} \\
  Device drivers & Everything that talks to hardware & \code{drivers/} \\
  Architecture & The parts that differ per CPU family & \code{arch/} \\
\end{booktable}

\noindent Userspace reaches all of it through system calls. Drivers reach hardware directly. The
layer that matters most to this course is the one that is not in the table: \textbf{the driver
model}, which is how a driver is matched to a device, and which lives partly in \file{drivers/base/}
and partly in every bus. Chapter~\ref{c10:ch} is about it.

\subsection{The tree, by the numbers}
\label{c3:app:a4}

The pinned source for this course is Linux 6.12.30. Measured:

\begin{narrowtable}{@{}lr@{}}
   & \thead{Count} \\
  \midrule
  Files & 86,660 \\
  C and header files & 59,998 \\
  Lines of C & 24,631,623 \\
  Lines of C and headers & 34,693,754 \\
\end{narrowtable}

\noindent Twenty-four million lines is not a number to be impressed by; it is a number to have a
strategy about. The strategy is that \textbf{you never read the kernel, you read one file in it},
and the tree is arranged so that you can find that file.

Where the bulk actually is:

\begin{booktable}{@{}lrlL@{}}
  \thead{Directory} & \thead{Files} & \thead{Size} & \thead{What is in it} \\
  \midrule
  \code{drivers/} & 34,970 & 1.1 GB & Device drivers. Two thirds of everything \\
  \code{arch/} & 17,417 & 153 MB & One subdirectory per CPU architecture \\
  \code{Documentation/} & 10,123 & 74 MB & Genuinely worth reading, and usually current \\
  \code{tools/} & 7,700 & 85 MB & Userspace tools shipped with the kernel \\
  \code{include/} & 6,206 & 56 MB & Headers, including the ones a module compiles against \\
  \code{sound/} & 2,781 & 52 MB & ALSA \\
  \code{fs/} & 2,480 & 51 MB & Filesystems, one subdirectory each \\
  \code{net/} & 1,992 & 38 MB & The network stack \\
  \code{kernel/} & 592 & 15 MB & Scheduler, timers, modules, tracing \\
  \code{mm/} & 196 & 5.8 MB & Memory management \\
\end{booktable}

\noindent Two observations are worth making.

\textbf{\file{drivers/} is two thirds of the kernel}, and within it, \file{drivers/gpu} alone is
553~MB, half of \file{drivers/} and a third of the whole tree. A course about embedded Linux will
never open that directory. This is why arm64 \code{defconfig} takes an hour to build, and why
\file{kernel/trim.config} exists.

\textbf{\file{arch/} is where portability lives, and where it stops.} \file{arch/arm64/} holds the
entry code, the page table format, the memory model and the syscall table. Everything above it is
written once. When something behaves differently on your target than on your laptop, \file{arch/} is
usually why.

\subsection{Finding things}
\label{c3:app:a5}

Four techniques, in the order you should reach for them.

\textbf{Know the layout.} A driver for hardware of type X is in \file{drivers/X/}. A filesystem is
in \file{fs/<name>/}. An architecture detail is in \file{arch/<arch>/}. This answers most questions
with no searching at all.

\textbf{Search for the string the user sees.} If \code{dmesg} printed something, that string is in
the tree:

\begin{shell}
grep -rn "loading out-of-tree module taints kernel" kernel/
\end{shell}

\noindent This is the single most effective kernel debugging technique there is, and it works
because the kernel prints in English and does not localise.

\textbf{Search for the interface, not the implementation.} To find who implements something, find
who registers it. \code{grep -rl "struct platform_driver" drivers/rtc/} finds every RTC platform
driver. Searching for \code{platform_driver_register} finds three, because the rest register
through the \code{module_platform_driver()} macro and never spell the function out.

\textbf{Read \file{Documentation/}.} It is unusually good and unusually current, and
\file{Documentation/driver-api/} in particular is the reference for most of what this course
teaches.

For reading rather than searching, Bootlin's
Elixir\footnote{\url{https://elixir.bootlin.com/linux/latest/source}} cross-references every
identifier in the tree, in a browser, without cloning anything.

\subsection{Reading one real driver}
\label{c3:app:a6}

The shape of a Linux driver is worth meeting before you write one.
\file{drivers/rtc/rtc-digicolor.c} is \textbf{224 lines}, drives the real-time clock on an SoC none
of us will ever own, and is a good first driver to read for exactly that reason: there is nothing
special about it. It is the ordinary shape.

Its skeleton, with the real line numbers from the pinned tree:

\begin{ccode}
static int __init dc_rtc_probe(struct platform_device *pdev)      /* 176: found my hardware */
{
        rtc->regs = devm_platform_ioremap_resource(pdev, 0);      /* 185: map registers   */
        rtc->rtc_dev = devm_rtc_allocate_device(&pdev->dev);      /* 189: join a framework */
        irq = platform_get_irq(pdev, 0);                          /* 193: which interrupt? */
        ret = devm_request_irq(&pdev->dev, irq, dc_rtc_irq, 0,    /* 196: take it          */
                               pdev->name, rtc);
        return devm_rtc_register_device(rtc->rtc_dev);            /* 205: go live          */
}

static const __maybe_unused struct of_device_id dc_dt_ids[] = {   /* 208: what I handle    */
        { .compatible = "cnxt,cx92755-rtc" },
        { }
};

static struct platform_driver dc_rtc_driver = {
        .driver = {
                .of_match_table = of_match_ptr(dc_dt_ids),        /* 217: how I am matched */
        },
};
module_platform_driver_probe(dc_rtc_driver, dc_rtc_probe);        /* 220                   */
\end{ccode}

\noindent \textbf{Every line of that maps onto a later chapter, in order.} \code{devm_} and
\code{ioremap} are Chapter~\ref{c6:ch}. \code{platform_get_irq} and \code{devm_request_irq} are
Chapter~\ref{c8:ch}. \code{of_match_table} and \code{compatible} are Chapter~\ref{c10:ch}.
\code{devm_rtc_register_device} is Chapter~\ref{c11:ch}. There is no \code{read()} and no
\file{/dev} node written by hand, and that absence is Chapter~\ref{c11:ch}'s argument.

Notice what is \emph{not} in it: no addresses, no interrupt numbers, no unwind path. The addresses
come from the device tree, and the \code{devm_} prefix means the cleanup is registered rather than
written. A driver written the way Chapter~\ref{c5:ch} teaches would be half again as long and would
have a ladder of \code{goto} labels at the bottom. Both of those disappearances are deliberate and
both are later chapters.

\textbf{The RTC actually on this target is a different one}, \file{drivers/rtc/rtc-pl031.c}, 470
lines. It is worth knowing that it is matched by a different mechanism: it is an ARM PrimeCell, so
it is identified by ID registers in the hardware itself rather than by a \code{compatible} string,
and it uses \code{amba_driver} rather than \code{platform_driver}. Bus-specific matching is
Chapter~\ref{c10:ch}'s subject; the reason to mention it now is that if you go looking at the driver
for the clock in your own machine, you will find it does not look like the skeleton above, and that
is not because one of them is wrong.

You can see it running:

\begin{console}
# dmesg | grep pl031
rtc-pl031 9010000.pl031: registered as rtc0
# grep pl031 /proc/interrupts
 15:          0          0  GIC-0  34 Level     rtc-pl031
\end{console}

\noindent Zero interrupts on both CPUs, because nothing has set an alarm. Chapter~\ref{c8:ch} is
where that column starts moving.

\subsection{Versions, and which tree you are on}
\label{c3:app:a7}

\begin{booktable}{@{}lLl@{}}
  \thead{Tree} & \thead{Meaning} & \thead{Lifetime} \\
  \midrule
  Mainline & Linus's tree. New features land here first & Continuous \\
  \code{-stable} & Backported fixes for the current release & \textasciitilde{}2 months \\
  \textbf{LTS} & A release designated long-term stable & 2 to 6 years \\
  Vendor & An SoC vendor's fork of some older release & Until they stop \\
\end{booktable}

\noindent \textbf{This course pins 6.12.30}, which is an LTS release. The \code{.30} is the stable
increment: 6.12 was the feature release, and 6.12.30 is that plus thirty rounds of backported fixes
and no new features. Stable releases take only fixes, judged against documented rules in
\file{Documentation/process/stable-kernel-rules.rst}, of which the important one is that a fix must
already be in mainline before it can be backported.

\textbf{The vendor tree is where embedded work actually happens, and it is the source of most of the
pain.} A typical SoC ships with a vendor kernel that is a fork of a release three or four years old,
with several thousand out-of-tree patches. It boots your board and mainline does not. What follows
from that:
\begin{itemize}
  \item A fix from mainline may not apply, because the code around it has moved.
  \item An upgrade is a port, not an update.
  \item Everything you write against vendor-specific APIs is stranded when you do upgrade.
\end{itemize}
The engineering response is to push whatever you can upstream and to prefer mainline interfaces even
when the vendor offers a shortcut. This course teaches mainline interfaces throughout, and that is
the reason.

\subsection{What ``24 MB'' is made of}
\label{c3:app:a8}

Chapter~\ref{c1:ch} measured the kernel this course builds at \textbf{about 24 MB} and left the
size unexplained. The explanation belongs here.

\code{vmlinux} is the linked kernel as an ELF object: \textbf{31,468,440 bytes}. \code{Image} is
that stripped of ELF structure and laid out for the bootloader to load: \textbf{24,943,104 bytes}.
The difference is symbol tables and ELF metadata that the boot process does not need.

Inside it, the largest contributors on this build:

\begin{narrowtable}{@{}lrl@{}}
  \thead{Section} & \thead{Size} & \thead{What it is} \\
  \midrule
  \code{.text} & 13,021,184 & Executable code \\
  \code{.rodata} & 4,369,606 & Constants, string literals, tables \\
  \code{.rela.dyn} & 3,341,232 & Relocations, applied once at boot and then freed \\
  \code{.init.*} & 1,029,519 & Code and data used once at boot and then freed \\
\end{narrowtable}

\noindent The last two rows are worth a moment, because you can watch them happen:

\begin{console}
[    5.001253] Freeing unused kernel memory: 4672K
\end{console}

\noindent Every function marked \code{__init} and every variable marked \code{__initdata} is placed
in a section that is discarded once boot is over, and on arm64 so is the table of relocations the
kernel applies to itself when it boots at a randomised address. On this build that is
\textbf{4,672 KB returned to the system}, the whole stretch from \code{__init_begin} to
\code{__init_end}, which is more than the entire BusyBox userland. Chapter~\ref{c4:ch} uses the same
annotations.

\textbf{And the reason it is 24 MB rather than 3 MB is configuration.} Nothing in the size is
mysterious: it is the sum of what was compiled in, and what was compiled in was chosen by a
\code{.config} with \textbf{5,725 lines} in it. \Secref{c3:app:b} is about that file, and the
Cross-check is about learning to predict what one line of it is worth.
